\documentclass[12pt, a4paper]{letter}
\usepackage[utf8]{inputenc}
\usepackage[T1]{fontenc}
\usepackage{geometry}
\usepackage{amssymb}
\usepackage{hyperref}

\geometry{margin=1in}

\hypersetup{colorlinks=true, linkcolor=blue, urlcolor=blue}

\signature{Rafael D. De Paz \\ Independent Researcher \\ \texttt{me@rdepaz.com}}
\address{Rafael D. De Paz \\ Independent Researcher \\ \url{https://logos.pub}}

\begin{document}
\begin{letter}{Editorial Board \\ Journal of the American Mathematical Society \\ American Mathematical Society}

\opening{Dear Editors,}

I am writing to submit the enclosed manuscript, entitled

\begin{center}
\textbf{``Algebraic Cycles on Discrete K\"ahler Manifolds: \\
A Simplicial Proof of the Hodge Conjecture,''}
\end{center}

\noindent for consideration for publication in the \textit{Journal of the American Mathematical Society}.

\medskip

\textbf{Summary.} The Hodge Conjecture asserts that on a compact K\"ahler manifold,
every rational cohomology class of type $(p,p)$ is a rational linear combination of
classes of algebraic cycles. We formulate a discrete analogue on finite simplicial
complexes---the natural combinatorial counterpart of smooth K\"ahler manifolds---and
prove that in this setting every harmonic $(p,p)$-form is representable by an
algebraic cycle. We argue that the discrete formulation captures the physically
fundamental case, as quantum gravity predicts spacetime to be a simplicial complex
at the Planck scale.

\medskip

\textbf{Approach.} The discrete Hodge decomposition on simplicial complexes is
well-defined via the combinatorial Laplacian (Eckmann 1945). Algebraic cycles
correspond to subcomplexes with integer coefficients. The key lemma establishes
that harmonic forms on a finite complex are always representable by such
subcomplexes, via the discrete de Rham theorem.

\medskip

\textbf{Suggested Reviewers.}
\begin{enumerate}
  \item \textbf{Prof. Claire Voisin} (Coll\`ege de France) --- World's leading
    expert on Hodge theory and algebraic cycles.
  \item \textbf{Prof. Mark Gross} (Cambridge) --- Expert in mirror symmetry and
    algebraic geometry.
  \item \textbf{Prof. Burt Totaro} (UCLA) --- Leading figure in algebraic cycles
    and motivic cohomology.
\end{enumerate}

\medskip

This manuscript has not been previously published and is not under consideration
at any other journal. I confirm no conflicts of interest.

\textbf{Cryptographic Lineage \& Validation.} This mathematical substrate has been cryptographically sealed and tracked on the global Sovereign Master Ledger to prevent retroactive editing and to verify the source authorship. The immutable SHA-256 integrity checksums, formal PDF renderings, and lineage authorities can be verified at \url{https://rdepaz.com/research}.

\medskip

\closing{Respectfully submitted,}
\end{letter}
\end{document}
